\documentclass[12pt]{article}
\usepackage[utf8]{inputenc}
\usepackage[T1]{fontenc}
\usepackage{hyperref}
\title{Barro Negro Tunebo Language Resources for K-12 Curriculum Development}
\author{Compiled by Claude Sonnet 5 (claude-sonnet-5) for Andr\'as Kornai}
\date{}
\begin{document}
\maketitle

\section{Introduction}
Barro Negro Tunebo (ISO 639-3 tbn; Glottocode barr1251), also called Eastern Tunebo or Yithkaya, is a Chibchan language of the U'wa people in Arauca and Casanare, Colombia. Speaker counts vary a great deal between sources: 300 (older Ethnologue, 1981), 500 (Global Recordings Network) and 3,800 (Joshua Project, Eastern Tunebo group). Little material was found.

\section{Primary Educational Resources}
None found.

\section{Linguistic and Reference Materials}
\begin{itemize}
\item Wikipedia, Uwa language (overview of the Tunebo varieties). \url{https://en.wikipedia.org/wiki/Uwa_language}
\item Chibchan documentation project, Tunebo (PDF, University of Missouri). Its content was not examined. \url{https://dice.missouri.edu/assets/docs/chibchan/Tunebo.pdf}
\item Glottolog reference: Pradilla Rueda, Elena (1985). \url{https://glottolog.org/resource/reference/id/31778}
\item Glottolog reference: Stoddart, D.R. (1962). \url{https://glottolog.org/resource/reference/id/159174}
\end{itemize}

\section{Religious and Community Resources}
\begin{itemize}
\item Joshua Project, Eastern Tunebo in Colombia (written status listed as unknown). \url{https://joshuaproject.net/people_groups/11708/CO}
\end{itemize}

\section{Digital and Audio Resources}
Global Recordings Network reports that it has no recordings in this language.
\begin{itemize}
\item Global Recordings Network, tbn page. \url{https://globalrecordings.net/en/language/tbn}
\item Endangered Languages Project, Tunebo entry. \url{https://endangeredlanguages.com/lang/3527}
\end{itemize}

\section{Curriculum Development Context}
The language is not known to be taught in schools, and no school programme, orthography or primer was found. Any curriculum would need to begin with community consultation with the U'wa and with locating descriptive work on this variety. Materials for the related Western Tunebo (tnb) and Central Tunebo (tuf) varieties should not be assumed to apply.
\end{document}